\documentclass[../readings.tex]{subfiles}

\begin{document}

\subsection{Chebyshev Polynomials}
\label{sub:chebyshev-theory}

\textBox{%
Chebyshev polynomials are the polynomial analogue of Fourier modes on $[-1,1]$. They oscillate evenly, stay bounded, and give node sets that control interpolation error. The central idea is minimax: among polynomials of a fixed degree, Chebyshev polynomials spread error as evenly as possible instead of allowing large boundary spikes.
}

\subsubsection{Failure of the Monomial Basis}

\textBox{%
Polynomial approximation has two separate failure modes. The first is numerical: the monomial basis $1,x,x^2,\dots$ leads to ill-conditioned Vandermonde matrices. The second is analytical: even exact interpolation at equally spaced nodes can oscillate badly near the interval endpoints.
}

\addDef{Interpolation Error Polynomial}{%
Given distinct interpolation nodes $x_1,\dots,x_n$, define
\begin{align*}
    \omega_n(x)=\prod_{j=1}^n (x-x_j).
\end{align*}
This polynomial is zero at every node. Its size between the nodes controls how much interpolation can amplify the derivative of the target function.
}
\label{def:cheb-error-polynomial}

\addThm{Interpolation Error Formula}{%
Let $p_{n-1}$ interpolate a function $f$ at $n$ distinct nodes $x_1,\dots,x_n$. If $f$ has $n$ continuous derivatives on $[-1,1]$, then for each $x\in[-1,1]$ there exists a point $\xi_x\in[-1,1]$ such that
\begin{align*}
    f(x)-p_{n-1}(x)
    =
    \frac{f^{(n)}(\xi_x)}{n!}\omega_n(x).
\end{align*}
Therefore,
\begin{align*}
    \|f-p_{n-1}\|_\infty
    \leq
    \frac{\|f^{(n)}\|_\infty}{n!}\|\omega_n\|_\infty.
\end{align*}
}
\label{thm:cheb-interpolation-error}

\addExcr{%
\begin{enumerate}
    \item For three nodes $x_1,x_2,x_3$, write out $\omega_3(x)$ and verify that it vanishes at each node.
    \item Use Thm~\ref{thm:cheb-interpolation-error} to explain why node choice matters even when $f^{(n)}$ is fixed.
    \item Build the degree-$14$ Vandermonde matrix on $15$ equally spaced nodes in $[-1,1]$. Compute $\kappa_2(\bV)$.
    \item Plot Runge's function $f(x)=1/(1+25x^2)$ and its degree-$10$ interpolant at equally spaced nodes. Compare the largest error in the middle of the interval with the largest error near the endpoints.
\end{enumerate}
}
\label{ex:cheb-monomial-failure}

\subsubsection{Chebyshev Polynomials ($T_k$)}

\addDef{Chebyshev Polynomial}{%
For $x\in[-1,1]$, the degree-$k$ Chebyshev polynomial is
\begin{align*}
    T_k(x)=\cos(k\arccos x).
\end{align*}
Equivalently, if $x=\cos\theta$, then
\begin{align*}
    T_k(\cos\theta)=\cos(k\theta).
\end{align*}
}
\label{def:cheb-polynomial}

\addThm{Basic Chebyshev Identities}{%
The polynomials from Def~\ref{def:cheb-polynomial} satisfy
\begin{align*}
    T_0(x)=1,\qquad T_1(x)=x,\qquad
    T_{k+1}(x)=2xT_k(x)-T_{k-1}(x).
\end{align*}
They are bounded by $1$ on $[-1,1]$:
\begin{align*}
    |T_k(x)|\leq 1.
\end{align*}
They are orthogonal with respect to the weight $(1-x^2)^{-1/2}$:
\begin{align*}
    \int_{-1}^1
    T_j(x)T_k(x)\frac{dx}{\sqrt{1-x^2}}
    =0,
    \qquad j\neq k.
\end{align*}
}
\label{thm:cheb-identities}

\addExcr{%
\begin{enumerate}
    \item Use $\cos((k+1)\theta)+\cos((k-1)\theta)=2\cos\theta\cos(k\theta)$ to prove the three-term recurrence in Thm~\ref{thm:cheb-identities}.
    \item Compute $T_2$, $T_3$, and $T_4$ from the recurrence.
    \item Prove $|T_k(x)|\leq 1$ on $[-1,1]$ directly from Def~\ref{def:cheb-polynomial}.
    \item Substitute $x=\cos\theta$ into the weighted integral and reduce the orthogonality statement to Fourier cosine orthogonality on $[0,\pi]$.
    \item Plot $T_0,\dots,T_5$. Mark the points where each polynomial reaches $\pm 1$.
\end{enumerate}
}
\label{ex:cheb-basic-identities}

\subsubsection{Zeros and Nodes}

\addThm{Zeros and Extrema}{%
The zeros of $T_n$ are
\begin{align*}
    x_j=\cos\left(\frac{(2j-1)\pi}{2n}\right),
    \qquad j=1,\dots,n.
\end{align*}
The extrema of $T_n$ occur at
\begin{align*}
    y_j=\cos\left(\frac{j\pi}{n}\right),
    \qquad j=0,\dots,n,
\end{align*}
and satisfy $T_n(y_j)=(-1)^j$.
}
\label{thm:cheb-zeros-extrema}

\addDef{Chebyshev Nodes}{%
The Chebyshev interpolation nodes are the zeros of $T_n$:
\begin{align*}
    x_j=\cos\left(\frac{(2j-1)\pi}{2n}\right),
    \qquad j=1,\dots,n.
\end{align*}
These nodes cluster near $\pm 1$. The clustering is not cosmetic; it reduces the size of the interpolation error polynomial from Def~\ref{def:cheb-error-polynomial}.
}
\label{def:cheb-nodes}

\addThm{Chebyshev Error Polynomial}{%
If $x_1,\dots,x_n$ are the Chebyshev nodes from Def~\ref{def:cheb-nodes}, then
\begin{align*}
    \omega_n(x)=\prod_{j=1}^n (x-x_j)
    =
    2^{1-n}T_n(x).
\end{align*}
Consequently,
\begin{align*}
    \|\omega_n\|_\infty=2^{1-n}.
\end{align*}
}
\label{thm:cheb-error-polynomial}

\addExcr{%
\begin{enumerate}
    \item Starting from $T_n(\cos\theta)=\cos(n\theta)$, solve $\cos(n\theta)=0$ and prove the zero formula in Thm~\ref{thm:cheb-zeros-extrema}.
    \item Solve $\cos(n\theta)=\pm 1$ and prove the extrema formula in Thm~\ref{thm:cheb-zeros-extrema}.
    \item Explain why the Chebyshev nodes cluster near $\pm 1$ by comparing the spacing of equally spaced $\theta$ values under $x=\cos\theta$.
    \item Show that $T_n$ has leading coefficient $2^{n-1}$ for $n\geq 1$.
    \item Use the previous step and the fact that $T_n$ has zeros $x_1,\dots,x_n$ to prove Thm~\ref{thm:cheb-error-polynomial}.
\end{enumerate}
}
\label{ex:cheb-nodes}

\subsubsection{The Minimax Property}

\addThm{Chebyshev Minimax Polynomial}{%
Among all monic polynomials $q$ of degree $n$ on $[-1,1]$,
\begin{align*}
    2^{1-n}T_n(x)
\end{align*}
has the smallest possible infinity norm:
\begin{align*}
    \|2^{1-n}T_n\|_\infty
    =
    2^{1-n}
    \leq
    \|q\|_\infty.
\end{align*}
}
\label{thm:cheb-minimax}

\addExcr{%
\begin{enumerate}
    \item Let $m(x)=2^{1-n}T_n(x)$. Use Thm~\ref{thm:cheb-identities} to show that $m$ is monic and $\|m\|_\infty=2^{1-n}$.
    \item Suppose another monic polynomial $q$ satisfies $\|q\|_\infty<2^{1-n}$. Define $r=q-m$. What is the degree of $r$?
    \item Evaluate the signs of $m$ at the $n+1$ extrema from Thm~\ref{thm:cheb-zeros-extrema}.
    \item Use the assumption $\|q\|_\infty<2^{1-n}$ to show that $r$ changes sign between each pair of adjacent extrema.
    \item Conclude that $r$ has at least $n$ roots, contradicting its degree. This proves Thm~\ref{thm:cheb-minimax}.
    \item Combine Thm~\ref{thm:cheb-interpolation-error} and Thm~\ref{thm:cheb-error-polynomial} to explain why Chebyshev nodes are the natural interpolation nodes.
\end{enumerate}
}
\label{ex:cheb-minimax}

\subsubsection{Chebyshev Series}

\addDef{Chebyshev Series}{%
A Chebyshev expansion writes a function on $[-1,1]$ as
\begin{align*}
    f(x)\sim \sum_{k=0}^{\infty} c_kT_k(x).
\end{align*}
Using $x=\cos\theta$, this is exactly a cosine series for the even function $f(\cos\theta)$.
}
\label{def:cheb-series}

\addThm{Coefficient Formula}{%
For $k\geq 1$,
\begin{align*}
    c_k=\frac{2}{\pi}\int_0^\pi f(\cos\theta)\cos(k\theta)\,d\theta,
\end{align*}
and
\begin{align*}
    c_0=\frac{1}{\pi}\int_0^\pi f(\cos\theta)\,d\theta.
\end{align*}
}
\label{thm:cheb-coefficients}

\addThm{Spectral Accuracy}{%
If $f$ is analytic on and near $[-1,1]$, then its Chebyshev coefficients decay geometrically:
\begin{align*}
    |c_k|\leq C\rho^{-k}
    \qquad \text{for some } \rho>1.
\end{align*}
The truncated series
\begin{align*}
    p_n(x)=\sum_{k=0}^n c_kT_k(x)
\end{align*}
therefore satisfies
\begin{align*}
    \|f-p_n\|_\infty=O(\rho^{-n}).
\end{align*}
This is called spectral accuracy.
}
\label{thm:spectral-accuracy}

\addRemark{%
\textbf{(Fourier connection)}\enspace Chebyshev methods are Fourier cosine methods after the change of variables $x=\cos\theta$. This is why fast cosine transforms can compute Chebyshev coefficients efficiently.
}

\addRemark{%
\textbf{(Interpolation vs.\ expansion)}\enspace Chebyshev interpolation chooses values at Chebyshev nodes. Chebyshev series choose coefficients in the Chebyshev basis. Both work well for the same reason: the basis oscillates evenly and avoids the endpoint instability of monomials.
}

\addExcr{%
\begin{enumerate}
    \item Use Def~\ref{def:cheb-series} to rewrite a Chebyshev expansion as a cosine series in $\theta$.
    \item Derive the coefficient formula in Thm~\ref{thm:cheb-coefficients} from Fourier cosine orthogonality.
    \item Compute degree-$n$ Chebyshev fits to $e^x$ using \texttt{np.polynomial.chebyshev.chebfit}. Plot $\|f-p_n\|_\infty$ versus $n$ on a semilog scale.
    \item Repeat the previous experiment for $f(x)=|x|$. Compare the coefficient decay with the analytic case in Thm~\ref{thm:spectral-accuracy}.
    \item Use \texttt{chebder} to differentiate a Chebyshev approximation to $f(x)=\sin(\pi x)$. Compare the derivative error with a finite-difference approximation using the same number of sample points.
\end{enumerate}
}
\label{ex:cheb-series}

\end{document}
